% !TEX root = ../main.tex
\phantomsection
\section*{Abstract}
\label{sec:abstract}

The financial system has generated broad-based growth impetus by allocating capital efficiently through lending. In this light, the advancement of DeFi is likewise expected to hinge critically on the degree of on-chain lending activity. However, DeFi platforms face constraints in executing collateral-backed loans due to difficulties in enforcing collateral assets, and they similarly confront limitations in offering unsecured credit, since traditional credit-scoring models that incorporate tangible-asset and employment information cannot be applied directly.


Existing PageRank-based methods, such as the Adaptive Weighted PageRank (AWP) algorithm, primarily rely on transaction volume and recency (time-decay) to assess wallet trustworthiness. While these approaches effectively capture eco- nomic flows and activity levels, they might incur substantial computational overhead by requiring exhaustive parsing of historical transactions.


This study proposes EndorseRank, a novel on-chain reputation metric leveraging ERC-20 approve and permit allowances as direct indicators of wallet endorsement. Each approve/permit action explicitly grants spending permissions to another wallet, inherently representing trust. By constructing a directed, weighted graph $G_{\text{approve}}$based solely on the most recent allowance amounts for each wallet pair, EndorseRank eliminates the need for historical time-decay adjustments and significantly reduces computational complexity.


The proposed EndorseRank algorithm utilizes PageRank methodology on this allowance-based graph. Our research aims to:


\begin{enumerate}
\item[1.] Develop a robust data pipeline to extract and manage the latest approve/permit allowances.
\item[2.] Construct and efficiently compute reputation scores from a sparse and endorsement- oriented graph.
\item[3.] Evaluate EndorseRank's predictive accuracy concerning on-chain financial risk (e.g., defaults on Aave) compared to established transaction-based methods, using metrics such as AUC, precision, and recall.
\item[4.] Conduct extensive computational benchmarks to demonstrate superior efficiency in runtime and memory usage.
\end{enumerate}
By addressing these objectives, EndorseRank seeks to provide a computation- ally efficient, semantically clear, and accurate reputation scoring mechanism suitable for real-world DeFi platforms, enabling better risk management and decision- making.


\textbf{\emph{Keywords--- }}\textbf{DeFi, On-Chain Reputation, On-Chain Financial Risk Prediction, Blockchain Data Analytics, Predictive Modeling, Wallet Endorsement Graph}
